\documentclass{article}
\usepackage{../math_template}

\author{Jonathan Kasongo}
\title{Imperial-Cambridge Math Contest 2024 Rd 2 --- P1/5}

\begin{document}
\maketitle

\begin{problem}{Imperial-Cambridge Math Contest 2024 Rd 2 --- P1/5}
\begin{enumerate}[label=(\alph*)]
\item Prove that there exist distinct positive integers $a_1, a_2,
\ldots, a_{2024}$ such that for each $i \in \{1,2,\ldots, 2024\}$,
$a_i$ divides $a_1 a_2 \cdots a_{i-1}a_{i+1} \cdots a_{2024} + 1$.
\item Prove that there exist distinct positive integers $b_1, b_2,
\ldots, b_{2024}$ such that for each $i \in \{1,2,\ldots, 2024\}$,
$b_i$ divides $b_1 b_2 \cdots b_{i-1}b_{i+1} \cdots b_{2024} + 2024$.
\end{enumerate}
\end{problem}

\begin{solution}{Solution}
\textbf{(a):} Let $a_1 = 3$, then define $a_n$ like so
$
a_i = (a_1 a_2 \cdots a_{i-1}) + f(i), \ i \geq 2
$
where $f : \mathbb{N} \rightarrow \{-1, 1\}$ is 1 if $i$ is even, and $-1$
if $i$ is odd. \vspace{0.2cm}

Now if $i$ is odd, then there are $2024 - i$ numbers in the set
$\{ a_{i+1}, \ldots, a_{2024} \}$. These numbers alternate between
$\pm 1$ modulo $a_i$ and since there are an odd number of terms the
product of elements in that set is $-1$ modulo $a_i$. Thus,

$$
a_1 a_2 \cdots a_{i-1}a_{i+1} \cdots a_{2024} + 1
\equiv
(a_i + 1)(-1) + 1 \equiv
0 \ \text{(mod } a_i).
$$

Similarly if $i$ is even, then the product of the elements in the set
$\{ a_{i+1}, \ldots, a_{2024} \}$ is 1 modulo $a_i$ because there is an
even number of elements in that set, and so

$$
a_1 a_2 \cdots a_{i-1}a_{i+1} \cdots a_{2024} + 1
\equiv
(a_i - 1)(1) + 1 \equiv
0 \ \text{(mod } a_i).
$$

To finish we just need to prove that all the $a_i$'s are distinct. We
notice that
\begin{align*}
a_i &\geq (a_1 a_2 \cdots a_{i-1}) - 1\\
&= (a_{i-1} \pm 1)a_{i-1} - 1\\
&\geq (a_{i-1} - 1)a_{i-1} - 1\\
&= (a_{i-1})^2 - a_{i-1} - 1
\end{align*}
But since $x^2 - x - 1 > x$ holds for $x > 1+\sqrt{2}$ and $a_1 = 3 >
1 + \sqrt{2}$, we have that $a_i \geq (a_{i-1})^2 - a_{i-1} - 1 > a_{i-1}$
and so $(a_n)$ is strictly increasing. This completes the proof.
$\square$\\

\textbf{(b):} We use a similar approach. Let $b_1 = 4048$, then
define $b_n$ like so $b_i = (b_1 b_2 \cdots b_{i-1}) + f(i), \ 2 \leq i \leq
2022$, where $f(i)$ has the same definition as in part (a).
Set $b_{2023} = 1$ and $b_{2024} = 2024$. \vspace{0.2cm}

If $i = 2023$ then we have nothing to prove since every positive integer
is divisible by 1. If $i = 2024$ then $a_1\cdots a_{2023} + 2024 \equiv
0 \ \text{(mod } b_{2024})$ because $b_1 = 4048 = 2\cdot 2024$. Now in what
proceeds assume $i \leq 2022$. \vspace{0.2cm}

If $i$ is odd then there are $2022 - i$ numbers in the set
$\{ b_{i+1}, \ldots, b_{2022} \}$. These numbers alternate between $\pm 1$
modulo $b_i$ and since there is an odd number of terms, the product of
elements in that set is $-1$ mod $b_i$. Thus we write

\[
\begin{aligned}
&\phantom{\equiv} b_1 b_2 \cdots b_{i-1} b_{i+1} \cdots b_{2024} + 2024\\
&\equiv (b_i + 1)(-1)\cdot b_{2023} \cdot b_{2024} + 2024\\
&\equiv
-2024 + 2024 \equiv 0 \ \text{(mod } b_i)
\end{aligned}
\]

Now if $i$ is even, then the product of the elements in the set
$\{ b_{i+1},\ldots, b_{2022} \}$ is 1 mod $b_i$ since there is an even
number of elements in that set. So we put

\[
\begin{aligned}
&\phantom{\equiv} b_1 b_2 \cdots b_{i-1} b_{i+1} \cdots b_{2024} + 2024\\
&\equiv (b_i - 1)(1)\cdot b_{2023} \cdot b_{2024} + 2024\\
&\equiv
-2024 + 2024 \equiv 0 \ \text{(mod } b_i)
\end{aligned}
\]

Now due to the arguement in (a) we know that the numbers $a_1, \cdots,
a_{2022}$ are all pairwise distinct and strictly greater than $4047$
and since $1$ and $2024$ are distinct and both strictly less than 4047,
the entire sequence is pairwise distinct as desired. This completes the
proof. $\blacksquare$

\end{solution}

\end{document}
